\documentclass[11pt]{article}

\usepackage[a4paper,margin=1in]{geometry}
\usepackage{amsmath,amssymb,amsthm,mathtools}
\usepackage{booktabs}
\usepackage{microtype}

\newtheorem{theorem}{Theorem}
\newtheorem{lemma}[theorem]{Lemma}
\newtheorem{proposition}[theorem]{Proposition}
\newtheorem{remark}[theorem]{Remark}

\newcommand{\dd}{\,\mathrm d}

\title{Atlas 43 at Low Density:\\A Balanced-Join Theorem}
\author{}
\date{}

\begin{document}

\maketitle

\section{Scope and result}

Let \(H\) be the house graph, equivalently the theta graph with path
lengths \(1,2,3\).  Its claimed Goodman-type lower bound is
\begin{equation}\label{eq:target}
 t(H,W)\geq B(p):=p(2p-1)\bigl(p^3+(1-p)^3\bigr),
 \qquad p=t(K_2,W)\geq\frac12.
\end{equation}
The general low-density problem remains open.  This note proves
\eqref{eq:target} on a large, infinite-dimensional class containing the
natural low-density extremal candidates.

Let \(X,Y\) be arbitrary graphons.  Their \emph{balanced complete join}
\(J(X,Y)\) is the graphon on two classes of measure \(1/2\), equal to \(X\)
and \(Y\) inside the two classes and equal to \(1\) across the classes.
No regularity, finite-step, or pointwise-smallness assumption is made on
\(X\) or \(Y\).

\begin{theorem}[Arbitrary interiors]\label{thm:arbitrary}
Let \(W=J(X,Y)\).  If
\[
 \frac12\leq p=t(K_2,W)
 \leq \frac{1+1/\sqrt5}{2}=0.7236067977\ldots,
\]
then \(W\) satisfies the Atlas \(43\) bound \eqref{eq:target}.
\end{theorem}

When the two interiors agree, a triangle--\(C_4\) tradeoff gives a modest
but useful extension.

\begin{theorem}[Identical interiors]\label{thm:symmetric}
Let \(W=J(X,X)\).  If
\[
 \frac12\leq p=t(K_2,W)\leq\frac{59}{80}=0.7375,
\]
then \(W\) satisfies the Atlas \(43\) bound \eqref{eq:target}.
\end{theorem}

These are structural low-side theorems, not global graphon theorems, so
they do not enlarge the green interval in the classification table.  Their
importance is that arbitrary high-amplitude perturbations inside the two
Turan classes are allowed.  Thus the unresolved obstruction to a global
low-density theorem is narrowed to incomplete or unbalanced cuts, rather
than the pure critical-addition face.

\section{The 32-colouring identity}

We use \(P_j\) for the path with \(j\) edges and \(Paw\) for a triangle
with a pendant edge.  Put
\[
 x=t(K_2,X),\qquad y=t(K_2,Y),\qquad r=\frac{x+y}{2}.
\]
The edge density of \(W=J(X,Y)\) is
\begin{equation}\label{eq:p-r}
 p=\frac12+\frac{x+y}{4}=\frac{1+r}{2},
 \qquad 2p-1=r.
\end{equation}

Assign each of the five vertices of \(H\) to one of the two join classes.
Edges crossing the classes contribute \(1\); the remaining edges induce a
graph in \(X\) or \(Y\).  Grouping the \(2^5=32\) assignments by their two
induced graphs gives the following exact identity:
\begin{align}
32t(H,J(X,Y))={}&t(H,X)+t(H,Y)\notag\\
&+2t(Paw,X)+2t(Paw,Y)\notag\\
&+y\,t(K_3,X)+x\,t(K_3,Y)\notag\\
&+t(C_4,X)+t(C_4,Y)\notag\\
&+2t(P_2,X)+2t(P_2,Y)\notag\\
&+2t(P_3,X)+2t(P_3,Y)\notag\\
&+4y\,t(P_2,X)+4x\,t(P_2,Y)\notag\\
&+2xy+2x+2y.\label{eq:32identity}
\end{align}
For reference, the multiplicities of the ordered induced pairs are
\begin{center}
\begin{tabular}{c|c}
ordered pair, together with its reversal&multiplicity in each orientation\\
\midrule
\((H,1)\)&\(1\)\\
\((Paw,1)\)&\(2\)\\
\((K_3,K_2)\)&\(1\)\\
\((C_4,1)\)&\(1\)\\
\((P_3,1)\)&\(2\)\\
\((P_2,1)\)&\(2\)\\
\((P_2,K_2)\)&\(4\)\\
\((K_2,1)\)&\(2\)\\
\((K_2,K_2)\)&\(2\) total
\end{tabular}
\end{center}
The entries sum to \(32\), and directly verify \eqref{eq:32identity}.

We shall use only the elementary path and \(4\)-cycle inequalities
\begin{equation}\label{eq:sidorenko-small}
 t(P_2,Z)\geq z^2,\qquad
 t(P_3,Z)\geq z^3,\qquad
 t(C_4,Z)\geq z^4,
 \qquad z=t(K_2,Z).
\end{equation}
The first is Jensen applied to the degree function, the second is the
Blakley--Roy path inequality, and the third follows twice from
Cauchy--Schwarz: if \(C_Z(u,v)=\int Z(u,w)Z(w,v)\dd w\), then
\[
 t(C_4,Z)=\int C_Z^2
 \geq\left(\int C_Z\right)^2
 =t(P_2,Z)^2\geq z^4.
\]

\section{Proof for arbitrary interiors}

\begin{proof}[Proof of Theorem~\ref{thm:arbitrary}]
All terms on the first three lines of \eqref{eq:32identity} are
nonnegative.  Dropping them and applying \eqref{eq:sidorenko-small} gives
\begin{align}
32t(H,W)\geq{}&x^4+y^4+2(x^2+y^2)+2(x^3+y^3)\notag\\
&+4(x^2y+xy^2)+2xy+2(x+y).\label{eq:join-lower}
\end{align}
Write
\[
 x=r+z,\qquad y=r-z.
\]
By \eqref{eq:p-r},
\[
 B(p)=\frac{r(1+r)(1+3r^2)}8.
\]
Expanding the right side of \eqref{eq:join-lower} and subtracting this
target yields
\begin{equation}\label{eq:general-gap}
 t(H,W)-B(p)
 \geq
 \frac{
 r^2(1-5r^2)
 +(6r^2+2r+1)z^2+z^4
 }{16}.
\end{equation}
If \(0\leq r\leq1/\sqrt5\), every term on the right is nonnegative.
Equation \eqref{eq:p-r} translates this condition exactly into the range
in the theorem.
\end{proof}

Notice that imbalance between the two internal edge densities contributes
the explicitly positive \(z^2\) and \(z^4\) terms in
\eqref{eq:general-gap}.  The symmetric case is therefore the limiting one
for this basic certificate.

\section{A triangle--\(C_4\) tradeoff}

The next lemma uses the fact that a graphon with few triangles must put
most of its two-step paths on nonedges.

\begin{lemma}\label{lem:tradeoff}
Let \(Z\) be a graphon of edge density \(s<1\), triangle density
\(\tau=t(K_3,Z)\), and \(4\)-cycle density \(c_4=t(C_4,Z)\).  Then
\begin{equation}\label{eq:tradeoff}
 c_4\geq s^4,
 \qquad
 c_4\geq\frac{(s^2-\tau)_+^2}{1-s}.
\end{equation}
\end{lemma}

\begin{proof}
The first inequality is \eqref{eq:sidorenko-small}.  Let
\(C_Z(u,v)=\int Z(u,w)Z(w,v)\dd w\) and
\(a=t(P_2,Z)=\int C_Z\geq s^2\).  Since
\[
 a-\tau=\int(1-Z)C_Z,
\]
Cauchy--Schwarz gives
\[
 (a-\tau)^2
 \leq\left(\int(1-Z)\right)
       \left(\int(1-Z)C_Z^2\right)
 \leq(1-s)c_4.
\]
If \(\tau<s^2\), use \(a-\tau\geq s^2-\tau\).  If
\(\tau\geq s^2\), the second assertion is vacuous.
\end{proof}

\begin{proof}[Proof of Theorem~\ref{thm:symmetric}]
Put \(s=t(K_2,X)\), so that \(p=(1+s)/2\).  Specializing
\eqref{eq:32identity} to \(X=Y\) gives
\begin{align}
16t(H,J(X,X))={}&t(H,X)+2t(Paw,X)+s\,t(K_3,X)+t(C_4,X)\notag\\
&+(2+4s)t(P_2,X)+2t(P_3,X)+s^2+2s.
\label{eq:symmetric-identity}
\end{align}
After dropping the first two nonnegative terms and applying the path
inequalities, it is enough to prove
\begin{equation}\label{eq:needed-tradeoff}
 t(C_4,X)+s\,t(K_3,X)\geq 6s^4-s^2.
\end{equation}

For \(s\leq1/\sqrt5\), the first inequality in
Lemma~\ref{lem:tradeoff} is already sufficient.  Assume henceforth that
\[
 \frac1{\sqrt5}\leq s\leq\frac{19}{40},
 \qquad q=1-s,
 \qquad \tau=t(K_3,X).
\]
Set
\[
 \tau_*:=s^2(1-\sqrt q).
\]
If \(\tau\geq\tau_*\), then the first inequality in
\eqref{eq:tradeoff} gives
\[
 t(C_4,X)+s\tau
 \geq s^4+s^3(1-\sqrt q).
\]
If \(\tau<\tau_*\), the second inequality gives
\[
 t(C_4,X)+s\tau
 \geq \frac{(s^2-\tau)^2}{q}+s\tau.
\]
The expression on the right is decreasing for
\(0\leq\tau\leq\tau_*\): its derivative at the right endpoint is
\[
 s-\frac{2s^2}{\sqrt q}\leq0,
\]
because \(4s^2\geq q\) throughout the present range.  Hence the same lower
bound holds in both cases:
\begin{equation}\label{eq:uniform-tradeoff}
 t(C_4,X)+s\tau
 \geq s^4+s^3(1-\sqrt{1-s}).
\end{equation}

It remains to check
\[
 h(s):=1-5s^2+s(1-\sqrt{1-s})\geq0.
\]
On the interval under consideration,
\[
 h'(s)=1-10s-\sqrt{1-s}+\frac{s}{2\sqrt{1-s}}<0;
\]
indeed \(1-10s<-3\), while the last two terms combine to
\((3s-2)/(2\sqrt{1-s})<0\).  Finally,
\[
 h\left(\frac{19}{40}\right)
 =\frac{111}{320}-\frac{19\sqrt{210}}{800}>0,
\]
where the strict inequality follows by squaring
\(555>38\sqrt{210}\).  Thus \eqref{eq:uniform-tradeoff} implies
\eqref{eq:needed-tradeoff}.  Since \(p=(1+s)/2\), the endpoint
\(s=19/40\) is \(p=59/80\).
\end{proof}

\section{Exact model families and the remaining obstruction}

If \(X\equiv r\), then \(J(X,X)\) is the regular two-block graphon with
cross value \(1\) and within-class value \(r=2p-1\).  A direct \(2\times2\)
calculation gives
\[
 t(H,J(X,X))
 =\frac{r(r+1)(r^4-r^3+5r^2+r+2)}{16}
\]
and the exact gap
\begin{equation}\label{eq:two-block-gap}
 t(H,J(X,X))-B(p)
 =\frac{r^2(1-r^2)^2}{16}\geq0.
\end{equation}
Thus this natural smoothing of the balanced bipartite graphon satisfies the
bound for the entire interval \(1/2\leq p\leq1\), although its gap is only
quadratic as \(p\downarrow1/2\).

For \(0\leq r\leq1/2\), take \(X\) itself to be a balanced two-block
graphon with cross value \(2r\) and diagonal value \(0\).  The resulting
four-block graphon has
\begin{equation}\label{eq:four-block-gap}
 t(H,W)-B(p)=\frac{r^2(1-4r^2)}{16}\geq0.
\end{equation}
It meets equality at \(r=0\) and \(r=1/2\), corresponding to the balanced
complete \(2\)- and \(4\)-partite graphons.  This family is numerically a
slightly stronger low-density competitor than the constant-interior family,
but still stays above the target.

The small-memory search script
\texttt{experiments/atlas43\_step\_search.py} optimizes both equal-part and
variable-part step graphons.  It is only a falsification tool, not a
certificate.  Runs with up to six parts found no violation down to
\(p=0.5001\); the smallest configurations were balanced joins or refinements
of them.  In particular, the retained certificate from the high-density
proof can be negative on these examples while the true gap in
\eqref{eq:two-block-gap} or \eqref{eq:four-block-gap} remains positive.

The next useful theorem would therefore control a graphon of the form
\[
 W=T_2+F-V,
\]
where \(F\geq0\) is arbitrary on the two diagonal cells and \(V\geq0\) is
the missing-edge kernel on the two cross cells.  Theorems
\ref{thm:arbitrary} and \ref{thm:symmetric} settle the full-amplitude face
\(V=0\) on substantial low-density intervals.  A global low-side result
now requires either:
\begin{enumerate}
 \item an absorption inequality showing that the positive pure-addition
 pair energy and the deletion derivative dominate all \(F\)-\(V\) mixed
 terms; or
 \item a symmetrization theorem showing that a low-density minimizer may be
 taken to have a complete balanced cut.
\end{enumerate}
Neither statement is proved here.

\end{document}
